\section{Experiment 2: MathVista Mathematical Reasoning Hallucination Detection}
\label{sec:mathvista}

\subsection{Direct Mode Baseline Results}

\Cref{tab:mathvista_direct} reports the MathVista evaluation results
for all four models in Direct mode.

\begin{table}[htbp]
  \centering
  \caption{MathVista Direct mode results.
    Best values are underlined.}
  \label{tab:mathvista_direct}
  \small
  \begin{tabular}{l cccc}
    \toprule
    Metric
      & GPT-5.4-mini & Gemini 2.5 Flash & Qwen3.5-35B & Qwen3-VL-235B \\
    \midrule
    Hallucination Rate & 45.1\% & 34.3\% & \underline{13.3\%} & 16.3\% \\
    Average Score      & 3.31   & 4.46   & \underline{5.00}   & 4.85 \\
    $n_{\text{hallucinated}}$ & 451 & 343 & \underline{133} & 163 \\
    \quad Faithfulness & 189 & 278 & \underline{98}  & 113 \\
    \quad Factuality   & 14  & 13  & \underline{5}   & 6 \\
    \quad Logical      & 248 & 52  & \underline{30}  & 44 \\
    None               & 549 & 657 & \underline{867} & 837 \\
    \bottomrule
  \end{tabular}
\end{table}

\begin{figure}[htbp]
  \centering
  \includegraphics[width=0.85\textwidth]{fig4_mathvista_composition.pdf}
  \caption{Hallucination type composition over hallucinated samples only
    in MathVista Direct mode.
    Each bar sums to the model's total hallucinated count $n$.}
  \label{fig:mathvista_pie}
\end{figure}

\subsubsection{Direct Mode Analysis}

\textbf{From POPE to MathVista: a richer error surface.}
On POPE the four-model HR range spans roughly 5 percentage points;
on MathVista it spans more than 30 points (13.3\%--45.1\%),
because the task simultaneously probes
visual perception, mathematical knowledge,
and multi-step reasoning.
From the worst performer GPT-5.4-mini (45.1\%)
to the best Qwen3.5-35B-A3B (13.3\%),
HR differs by a factor of 3.4---
far exceeding the cross-model gap on POPE.
The capacity--calibration view from \Cref{sec:pope_cross_model_comparison}
remains useful as a starting point,
but each model also exposes a distinct dominant error type
that cannot be read off POPE alone.

\textbf{Distinct dominant error types.}
GPT-5.4-mini is dominated by logical hallucinations
(55.0\%, 248 out of 451).
In Direct mode this model typically outputs only an option letter
(e.g., ``B'') with a median response length of just one word,
providing no reasoning process for the judge to assess.
Many incorrect answers are thus directly classified as logical hallucinations,
reflecting both weak mathematical reasoning
and an ultra-short response style
that offers no reasoning chain to earn partial credit.

Gemini 2.5 Flash is dominated by faithfulness hallucinations
(81.0\%, 278 out of 343).
This model produces highly detailed responses
(median 242 words) with extensive visual descriptions,
but the descriptions contain visual misreadings,
such as incorrectly reporting chart values, colors, or labels.
Gemini's weakness lies not in reasoning (only 52 logical cases)
but in imprecise visual perception:
describing too much leads to getting too much wrong.

Both Qwen models achieve hallucination rates
less than one-third of GPT-5.4-mini.
Qwen3.5-35B-A3B (13.3\%) and Qwen3-VL-235B (16.3\%)
also show a more balanced distribution
across the three hallucination types,
with no obvious single weakness.
This contrasts with POPE results,
where the smaller Qwen3.5 (3\,B activated) trails
the 235\,B model slightly,
but on MathVista's mathematical reasoning it leads by a small margin,
suggesting higher efficiency in mathematical reasoning
despite its smaller activated parameter count.

\subsection{CoT Mode Results and Direct Comparison}

\Cref{tab:mathvista_cot_comparison} compares
the core metrics between Direct and CoT modes,
including overall performance changes and per-type hallucination breakdowns.

\begin{table}[htbp]
  \centering
  \caption{MathVista Direct vs.\ CoT comparison with per-type hallucination changes.}
  \label{tab:mathvista_cot_comparison}
  \small
  \setlength{\tabcolsep}{3pt}
  \begin{tabular}{l cc c cc c}
    \toprule
    \multicolumn{7}{c}{\textbf{Overall Performance Comparison}} \\
    \midrule
    Model & Direct HR & CoT HR & $\Delta$HR
      & Direct Avg & CoT Avg & $\Delta$Avg \\
    \midrule
    GPT-5.4-mini     & 45.1\% & 34.4\% & $\downarrow$10.7\,pp & 3.31 & 4.35 & +1.04 \\
    Gemini 2.5 Flash & 34.3\% & 27.6\% & $\downarrow$6.7\,pp  & 4.46 & 4.81 & +0.35 \\
    Qwen3.5-35B      & 13.3\% & 15.5\% & $\uparrow$2.2\,pp    & 5.00 & 5.34 & +0.34 \\
    Qwen3-VL-235B    & 16.3\% & 24.0\% & $\uparrow$7.7\,pp    & 4.85 & 4.94 & +0.09 \\
    \midrule
    \multicolumn{7}{c}{\textbf{Per-Type Hallucination Changes}} \\
    \midrule
    Model & \multicolumn{2}{c}{$\Delta$ Faithfulness} & $\Delta$ Factuality
      & \multicolumn{2}{c}{$\Delta$ Logical} & Net improvement \\
    \midrule
    GPT-5.4-mini     & \multicolumn{2}{c}{+8}   & $-$10 & \multicolumn{2}{c}{$-$105} & +107 \\
    Gemini 2.5 Flash & \multicolumn{2}{c}{$-$66} & $-$8  & \multicolumn{2}{c}{+7}     & +67  \\
    Qwen3.5-35B      & \multicolumn{2}{c}{+26}  & $-$3  & \multicolumn{2}{c}{$-$1}   & $-$22 \\
    Qwen3-VL-235B    & \multicolumn{2}{c}{+41}  & 0     & \multicolumn{2}{c}{+36}    & $-$77 \\
    \bottomrule
  \end{tabular}
  \vspace{0.5em}

  \small
  \textit{Note:} Net improvement denotes the net change in sample count (positive = improvement, negative = degradation).
\end{table}

\begin{figure}[htbp]
  \centering
  \includegraphics[width=0.8\textwidth]{fig2_mathvista_direct_cot_bar.pdf}
  \caption{Comparison of hallucination rates
    between Direct and CoT modes.
    $\Delta$ values indicate the change.}
  \label{fig:mathvista_cot_bar}
\end{figure}

\subsection{CoT Effect Analysis}

\subsubsection{Overview}

\textbf{Baseline-dependent CoT effects.}
CoT effects on MathVista divide cleanly along a single axis:
the Direct-mode baseline.
Models with high Direct HR
(GPT-5.4-mini at 45.1\%, Gemini at 34.3\%)
improve substantially under CoT,
whereas models with low Direct HR
(Qwen3.5 at 13.3\%, Qwen3-VL at 16.3\%)
degrade.
\Cref{tab:mathvista_cot_comparison,fig:mathvista_cot_bar}
quantify this pattern.
The remainder of this subsection isolates
the mechanisms responsible for each direction
and identifies why CoT helps logical hallucinations
but exacerbates faithfulness hallucinations.

\subsubsection{CoT Beneficiaries: Weaker Baselines}

CoT significantly improves weaker-baseline models.
GPT-5.4-mini is the largest beneficiary:
HR drops by 10.7 percentage points
and average score increases by 1.04,
both the largest improvements among all models.
The improvement comes primarily
from a sharp reduction in logical hallucinations
(248 $\to$ 143, $-$42.3\%),
as shown in \Cref{tab:mathvista_cot_comparison}.
In Direct mode this model outputs only option letters
(median one word),
exposing incorrect answers directly to the judge
without any reasoning evidence.
CoT forces the model to organize its mathematical derivation
step by step,
correcting many errors that a reasoning chain can catch.
Sample-level analysis shows that 209 Direct-mode hallucinations
are corrected under CoT,
while only 102 cases exhibit reverse degradation,
yielding a net improvement of 107 samples.

Gemini 2.5 Flash improves through a different pathway:
the benefit comes primarily from reduced faithfulness hallucinations
(278 $\to$ 212, $-$23.7\%).
CoT requires the model to explicitly cite visual evidence
during reasoning,
forcing more careful examination of chart content
and reducing visual misreadings.
Notably, Gemini's Direct-mode responses are already highly detailed
(median 242 words),
and its CoT responses are actually slightly shorter
(median 228 words),
indicating that CoT does not force longer output
but rather improves accuracy
by constraining reasoning to image-visible evidence.

\subsubsection{CoT Risks: Stronger Baselines}

CoT introduces additional risk for stronger-baseline models.
Both Qwen models exhibit a polarization pattern:
6-score and 2-score bins both grow
while intermediate 4--5 scores collapse,
indicating that CoT pushes responses toward the extremes
of either perfect correctness or detailed incorrectness.

Qwen3.5-35B-A3B exhibits a pattern
in which HR rises by 2.2 percentage points
while average score also rises from 5.00 to 5.34.
Score polarization explains both moves:
6-score samples grow from 488 to 810
(lifting the mean),
and 2-score samples grow from 125 to 154
(detailed but incorrect reasoning is strictly penalized,
lifting HR).
Since this model already possesses strong mathematical reasoning
in Direct mode (only 30 logical cases),
CoT yields little marginal benefit
and instead exposes blind spots in visual perception
by forcing explicit visual description,
causing faithfulness hallucinations to rise
from 98 to 124 (+26 cases, +26.5\%).

Qwen3-VL-235B shows a more pronounced pattern:
HR jumps from 16.3\% to 24.0\% (+7.7pp),
with a net degradation of 77 samples.
Of these, 116 Direct-mode correct samples
produce hallucinations under CoT,
while only 39 improve in the reverse direction.
The polarization effect is equally pronounced:
6-score samples surge from 452 to 723,
but 2-score samples increase from 142 to 194
and 0-score cases appear (0 $\to$ 8),
while intermediate 4--5 scores nearly vanish
(93 $\to$ 6 and 291 $\to$ 30 respectively).
Faithfulness hallucinations increase from 113 to 154 (+41 cases, +36.3\%),
indicating that forcing explicit visual description
systematically exposes perceptual weaknesses.

\subsubsection{Underlying Mechanisms}

This phenomenon aligns with recent findings.
\citet{liuMoreThinkingLess2025} show through attention analysis
that as reasoning chains grow longer,
the model's attention to visual input weakens (attention drift),
gradually shifting from image content to linguistic priors,
leading to ``more thinking, less seeing.''
\citet{kanchetiChainThoughtDegrades2025} report,
after evaluating 17 models,
that CoT consistently degrades visual spatial reasoning,
with models ``hallucinating visual details from textual priors.''
The substantial increase in faithfulness hallucinations for both Qwen models
is consistent with both observations.

CoT introduces hallucinations through two complementary mechanisms.

\begin{description}
  \item[Exposure effect.] Visual perception errors that remain hidden in Direct mode
    are explicitly described in the CoT reasoning chain
    and captured by the MLLM Judge.
    Even when the final answer is correct,
    visual misreadings in intermediate steps lower the score.
    This mechanism explains why stronger-baseline models
    see increased HR despite high final-answer accuracy.

  \item[Commitment error.] CoT forces the model to make explicit visual or mathematical assertions
    early in the reasoning process.
    Once an intermediate step goes wrong,
    the model continues reasoning around the incorrect premise,
    producing ``informative but wrong'' responses (score 2).
    The Multimodal Hallucination Snowballing framework
    proposed by \citet{zhongInvestigatingMitigatingMultimodal2024}
    directly corresponds to this mechanism:
    once the model produces a hallucination early in reasoning,
    it elaborates coherently around the erroneous premise,
    causing the error to propagate and amplify along the chain.
\end{description}

These two mechanisms operate together:
the exposure effect reveals latent perceptual errors,
while the commitment error transforms them into cascading reasoning failures.
Together they explain why CoT can simultaneously improve weaker models
(by catching logical errors through structured reasoning)
and degrade stronger models
(by forcing explicit articulation of visual misperceptions
that would otherwise remain implicit and harmless).

\subsection{Failure Case Analysis}

This section presents two typical failure cases,
each corresponding to one of the two core CoT side-effect mechanisms.

\subsubsection{Case 1: CoT Commitment Error Leading to Logical Hallucination
  (GPT-5.4-mini, PID=30)}

\Cref{tab:mathvista_case1_info} presents the sample information,
and \Cref{fig:mathvista_case1} shows the accompanying image.
The question is a circle geometry problem:
two chords $AB$ and $CD$ intersect at $E$,
with $\angle D = 35^\circ$ and $\angle AEC = 105^\circ$;
the task is to find $\angle C$.
The correct answer is $70^\circ$ (option B).

\begin{table}[htbp]
  \centering
  \caption{MathVista Case 1: sample information.}
  \label{tab:mathvista_case1_info}
  \small
  \begin{tabular}{ll}
    \toprule
    Field & Content \\
    \midrule
    Question type & Geometry reasoning (multiple choice) \\
    Question & Two chords $AB$, $CD$ intersect at $E$; $\angle D = 35^\circ$. Find $\angle C$. \\
    Ground Truth & $70^\circ$ (B) \\
    \bottomrule
  \end{tabular}
\end{table}

\begin{figure}[htbp]
  \centering
  \includegraphics[width=0.25\textwidth]{mathvista_pid30.jpg}
  \caption{MathVista Case 1: intersecting chords in a circle.}
  \label{fig:mathvista_case1}
\end{figure}

\begin{table}[H]
  \centering
  \caption{MathVista Case 1: responses under two modes.}
  \label{tab:mathvista_case1_answers}
  \small
  \begin{tabular}{llcc}
    \toprule
    Mode & Response & Score & Type \\
    \midrule
    Direct & B ($70^\circ$) & 5 & none \\
    CoT    & Flawed reasoning, arrives at $80^\circ$ (C) & 1 & logical \\
    \bottomrule
  \end{tabular}
\end{table}

In Direct mode the model outputs the correct option B ($70^\circ$)
without showing any reasoning, and the judge assigns score 5.
In CoT mode the model attempts step-by-step reasoning
but commits to an incorrect premise in the first step,
treating $\angle E = 105^\circ$ as a given condition
and applying the triangle angle-sum formula
to obtain $\angle C = 180^\circ - 105^\circ - 35^\circ = 40^\circ$.
Finding that $40^\circ$ is not among the choices,
the model then fabricates a supplementary-angle argument
to produce $80^\circ$ and selects option C.
This exhibits the snowballing hallucination pattern
described by \citet{zhongInvestigatingMitigatingMultimodal2024}:
once an incorrect premise is locked in,
subsequent steps build upon it,
and each step appears plausible in isolation
while the entire chain rests on a false initial assumption.

The correct approach uses the inscribed angle theorem:
$\angle AEC$ is the vertical angle of $\angle DEB$,
and $\angle C = \angle AEC - \angle D = 105^\circ - 35^\circ = 70^\circ$.
In Direct mode the model retrieves this result directly,
but CoT forces it to articulate intermediate steps
where it selects an incorrect decomposition strategy.
This demonstrates that CoT can convert a correct intuition
into a wrong answer when the model lacks
the precise geometric knowledge to sustain a valid reasoning chain.

\subsubsection{Case 2: CoT Exposure of Visual Perception Error
  Leading to Faithfulness Hallucination
  (Qwen3-VL-235B, PID=15)}

\Cref{tab:mathvista_case2_info} presents the sample information,
and \Cref{fig:mathvista_case2} shows the accompanying image.
The question asks which organism would be most affected
if algae were eliminated from a food web.
The correct answer is Common water flea (option B).

\begin{table}[htbp]
  \centering
  \caption{MathVista Case 2: sample information.}
  \label{tab:mathvista_case2_info}
  \small
  \begin{tabular}{ll}
    \toprule
    Field & Content \\
    \midrule
    Question type & Biology food web reasoning (multiple choice) \\
    Question & Which organism will be most affected if algae is eliminated? \\
    Ground Truth & Common water flea (B) \\
    \bottomrule
  \end{tabular}
\end{table}

\begin{figure}[htbp]
  \centering
  \includegraphics[width=0.35\textwidth]{mathvista_pid15.jpg}
  \caption{MathVista Case 2: food web diagram.}
  \label{fig:mathvista_case2}
\end{figure}

\begin{table}[H]
  \centering
  \caption{MathVista Case 2: responses under two modes.}
  \label{tab:mathvista_case2_answers}
  \small
  \begin{tabular}{llcc}
    \toprule
    Mode & Response & Score & Type \\
    \midrule
    Direct & B (correct answer) & 5 & none \\
    CoT    & B (correct, but reasoning contains visual error) & 2 & faithfulness \\
    \bottomrule
  \end{tabular}
\end{table}

In Direct mode the model outputs the correct answer
and receives score 5 (no hallucination).
In CoT mode the model provides a detailed food web analysis
but incorrectly claims that ``Algae has arrows pointing to Tilapia''
when in fact Tilapia obtains energy
from Common water flea and Water plant, not directly from Algae.
This is a typical exposure effect:
in Direct mode the model may reach the correct answer
through some shortcut (linguistic prior or partial visual perception)
while its visual perception is not fully accurate.
CoT forces the model to explicitly describe
the arrow relationships in the image,
exposing the visual misreading to the MLLM Judge
and causing the score to drop from 5 to 2.

This case reveals an important phenomenon in CoT evaluation:
a correct answer does not imply correct reasoning.
Under CoT the MLLM Judge evaluates both the reasoning process
and the final answer;
even when the answer is correct,
visual inconsistencies in the reasoning lower the score.
This explains why Qwen3-VL-235B
sees faithfulness hallucinations increase from 113 to 154 under CoT:
a substantial portion are cases
where the answer is correct but the reasoning contains visual errors,
errors that are hidden in Direct mode but exposed under CoT.
